\section*{अध्याय \ref{ch:b3:organometallic-bonding} --- कार्बधात्विक रसायन: आबंधन और लिगैंड}

\begin{solution}{exo:b3:organometallic-bonding:1}
$\ce{Fe(CO)5}$: $8 + 10 = 18$। $\ce{Mn2(CO)10}$: प्रति Mn $7 + 10 + 1 = 18$। $\ce{CpMo(CO)3H}$: $6 + 5 + 6 + 1 =
18$। $\ce{[PtCl3(C2H4)]-}$: $10 + 3 + 2 + 1 = 16$। $\ce{Cr($\eta^6$-C6H6)(CO)3}$: $6 + 6 + 6 = 18$। $\ce{Cp2ZrCl2}$:
$4 + 10 + 2 = 16$।
\end{solution}

\begin{solution}{exo:b3:organometallic-bonding:2}
Fe(0) $d^8$; Mn(0) $d^7$; Mo(II) $d^4$; Pt(II) $d^8$; Cr(0) $d^6$; Zr(IV) $d^0$।
\end{solution}

\begin{solution}{exo:b3:organometallic-bonding:3}
$\ce{[Mn(CO)6]+} > \ce{Cr(CO)6} > \ce{[V(CO)6]-}$: धातु जितनी अधिक इलेक्ट्रॉन-समृद्ध होती है, वह उतना ही अधिक
$\pi^*$ में पश्च-दान करती है और C--O आबंध उतने ही दुर्बल होते हैं।
\end{solution}

\begin{solution}{exo:b3:organometallic-bonding:4}
एथेन ($\ce{Mn(CO)5}$ और $\ce{CpFe(CO)2}$, \ce{CH3} के साथ \omterm{def:b3:organometallic-bonding:isolobal}{समपालि} हैं); फिर एथेन; टेट्राहेड्रेन
$\ce{(CH)4}$, जिसमें तीन CH को $\ce{Co(CO)3}$ से प्रतिस्थापित किया गया है।
\end{solution}

\begin{solution}{exo:b3:organometallic-bonding:5}
$fac$ ($C_{3v}$): दो (A$_1$ + E)। $mer$ ($C_{2v}$): तीन (2A$_1$ + B$_1$)।
\end{solution}

\begin{solution}{exo:b3:organometallic-bonding:6}
भारी फ़ॉस्फ़ीन धातु के आसपास भीड़ कर देता है: वियोजन तनाव को हल्का करता है, अतः वह
तेज़ होता है, और निम्न-उपसहसंयोजी मध्यवर्ती अनुकूल होते हैं।
\end{solution}

\begin{solution}{exo:b3:organometallic-bonding:7}
क्रोमियम संकुल एक \omterm{def:b3:organometallic-bonding:carbene}{फ़िशर कार्बीन} है (OMe प्रतिस्थापी, Cr(0), CO लिगैंड): कोई
ऐमीन इलेक्ट्रॉनरागी \omterm{def:b3:organometallic-bonding:carbene}{कार्बीन} कार्बन पर आक्रमण करके OMe को प्रतिस्थापित करती है। टैंटलम
ऐल्किलिडीन एक \omterm{def:b3:organometallic-bonding:carbene}{श्रॉक कार्बीन} है: उसका नाभिकरागी कार्बन किसी कीटोन के
कार्बोनिल पर आक्रमण करता है, जिससे एक ऐल्कीन और एक Ta=O इकाई बनती है, जैसे विटिग अभिक्रिया में।
\end{solution}

\begin{solution}{exo:b3:organometallic-bonding:8}
सांतरित: कोई $\delta$ अतिव्यापन नहीं, दोनों $\delta$ इलेक्ट्रॉन अनाबंधी हैं: आबंध कोटि 3। $\ce{[Re2Cl8]^{4-}}$:
$\sigma^2\pi^4\delta^2\delta^{\ast2}$, आबंध कोटि 3।
\end{solution}

\begin{solution}{exo:b3:organometallic-bonding:9}
छोटी M$\cdots$H दूरी (विवर्तन, सर्वोत्तम न्यूट्रॉन); कम C--H तनन
तरंग-संख्या; उच्च क्षेत्र पर घटे हुए एक-आबंध C--H युग्मन वाला NMR संकेत।
\end{solution}

\begin{solution}{exo:b3:organometallic-bonding:10}
16: वर्ग-समतलीय $d^8$ ($\ce{[PtCl4]^{2-}}$, विल्किन्सन उत्प्रेरक), जिनका 18वें इलेक्ट्रॉन का कक्षक
बहुत ऊँचा है; $\ce{Cp2ZrCl2}$ जैसे आरंभिक-धातु \omterm{def:b3:organometallic-bonding:metallocene}{मेटैलोसीन}, जो और
लिगैंडों के लिए बहुत भीड़ भरे हैं। 17: $\ce{V(CO)6}$, जो त्रिविम कारणों से द्वितय नहीं बना सकता। 19: कोबाल्टोसीन,
जिसका अतिरिक्त इलेक्ट्रॉन एक प्रतिआबंधी कक्षक में बैठता है (और आसानी से निकल जाता है)।
\end{solution}

\begin{solution}{exo:b3:organometallic-bonding:11}
CO एक प्रबल $\pi$ ग्राही है: यह $t_{2g}$ स्तर को नीचे करता है और $\Delta_{\mathrm o}$ को बड़ा, युग्मन ऊर्जा से
कहीं ऊपर, बनाता है: $t_{2g}^6$, प्रतिचुंबकीय। $d$--$d$ संक्रमण, $\Delta_{\mathrm o}$ पर, पराबैंगनी में पड़ते हैं,
और कोई दृश्य प्रकाश अवशोषित नहीं होता।
\end{solution}

\begin{solution}{exo:b3:organometallic-bonding:12}
$\pi$ से दान और $\pi^*$ में पश्च-दान, दोनों C--C आबंध कोटि घटाते हैं। ज़ाइस
लवण (Pt(II), मध्यम पश्च-दान) में आबंध थोड़ा लंबा होता है; एक निम्न-संयोजी,
इलेक्ट्रॉन-समृद्ध धातु प्रबल पश्च-दान करती है और संकुल
\omterm{def:b3:organometallic-bonding:dcd}{धातुसाइक्लोप्रोपेन} सीमा के निकट पहुँचता है, जिसका C--C आबंध एकल आबंध के निकट होता है।
\end{solution}

\begin{solution}{pb:b3:organometallic-bonding:1}
\textbf{1.} $\ce{Fe^{2+}}$ ($d^6$) $+ 2 \times 6 = 18$।
\textbf{2.} $8 + 2 \times 5 = 18$।
\textbf{3.} $9 + 10 = 19$।
\textbf{4.} $10 + 10 = 20$।
\textbf{5.} $18 - 1 = 17$।
\textbf{6.} Fe(II) $d^6$; Co(II) $d^7$; Ni(II) $d^8$; Fe(III) $d^5$।
\textbf{7.} $a_{1g}$ ($d_{z^2}$), $e_{1g}$ ($d_{xz}$, $d_{yz}$), $e_{2g}$ ($d_{xy}$, $d_{x^2-y^2}$)।
\textbf{8.} $e_{1g}$ युग्म भरे वलय $e_{1g}$ संयोजन से प्रबलता से अतिव्यापन करता है: आबंधी
कक्षक अधिकतर वलयों पर, प्रतिआबंधी $e_{1g}^\ast$ अधिकतर धातु पर, ऊर्जा में
ऊँचे।
\textbf{9.} $e_{2g} \approx a_{1g} < e_{1g}^\ast$।
\textbf{10.} $e_{2g}^4a_{1g}^2$: कोई अयुग्मित इलेक्ट्रॉन नहीं।
\textbf{11.} $e_{1g}^\ast$ में एक इलेक्ट्रॉन: एक अयुग्मित इलेक्ट्रॉन, $1.73\,\mu_B$।
\textbf{12.} दोनों $e_{1g}^\ast$ कक्षकों में दो इलेक्ट्रॉन, समांतर: दो अयुग्मित, $\sqrt{8} =
2.83\,\mu_B$।
\textbf{13.} $e_{2g}^3a_{1g}^2$: एक अयुग्मित इलेक्ट्रॉन। $e_{2g}^3$ छिद्र एक कक्षकीय रूप से अपभ्रष्ट
(E) अवस्था देता है, \omterm{def:b3:complex-spectra-magnetism:moment}{कक्षकीय योगदान} के साथ।
\textbf{14.} उसका 19वाँ इलेक्ट्रॉन एक प्रतिआबंधी कक्षक में है और आसानी से हट जाता है,
जिससे 18-इलेक्ट्रॉन कोबाल्टोसीनियम धनायन बनता है।
\textbf{15.} हटाया गया इलेक्ट्रॉन लगभग अनाबंधी कक्षक से आता है:
संरचना लगभग नहीं बदलती (छोटी \omterm{def:b3:complex-mechanisms:reorganisation}{पुनर्गठन ऊर्जा}, \cref{ch:b3:complex-mechanisms})।
\textbf{16.} युग्म अधिकांश विलायकों में तीव्र और उत्क्रमणीय है, यौगिक
विलेय और स्थायी है, और इसका विभव विलायक पर बहुत कम निर्भर करता है क्योंकि
आवेश एक बड़े आयन पर फैला होता है।
\textbf{17.} 20 इलेक्ट्रॉनों के साथ, जिनमें दो प्रतिआबंधी हैं, यह उन्हें खोकर अभिक्रिया करता है: यह
वलय के खिसकने या वलय के निकलने के साथ लिगैंड जोड़ता है, या ऑक्सीकृत होता है।
\textbf{18.} $\ce{CpFe(CO)2}$: $8 + 5 + 4 = 17$ इलेक्ट्रॉन, 18 से एक कम, एक अग्रांत कक्षक: \ce{CH3} के साथ
\omterm{def:b3:organometallic-bonding:isolobal}{समपालि}। द्वितय $\ce{[CpFe(CO)2]2}$ ``एथेन'' है, एक Fe--Fe आबंध के साथ (व्यवहार में दो CO
लिगैंड इसे सेतुबद्ध करते हैं)।
\textbf{19.} $\ce{CpFe(CO)2-Mn(CO)5}$, एक Fe--Mn आबंध के साथ।
\textbf{20.} $\ce{PPh3}$, CO की तुलना में बेहतर दाता और कमज़ोर $\pi$ ग्राही है: धातु शेष CO में
अधिक पश्च-दान करती है, जिनके तनन गिर जाते हैं।
\textbf{21.} दो (A$_1$ + E)।
\textbf{22.} $\eta^3$ तक वलय का खिसकना दो इलेक्ट्रॉनों के बराबर उपसहसंयोजन मुक्त करता है: प्रवेशी
लिगैंड धातु के 18 इलेक्ट्रॉन से आगे गए बिना सहयोजी रूप से बंध सकता है।
\textbf{23.} \textbf{$\mu_{\text{केवल चक्रण}}(\text{निकेलोसीन}) = \sqrt{2 \times 4} = 2.83\,\mu_B$।}
\end{solution}
